\documentclass[12pt,a4paper]{article}
\usepackage[utf8]{inputenc}
\usepackage[T5]{fontenc}
\usepackage[margin=2.2cm]{geometry}
\usepackage{xcolor}
\usepackage{mdframed}
\usepackage{enumitem}
\usepackage{parskip}
\usepackage{lmodern}
\usepackage{booktabs}
\usepackage{amsmath}
\usepackage{amssymb}
\usepackage{array}
\usepackage{tabularx}

% ── Màu sắc ─────────────────────────────────────────────────
\definecolor{darkblue}{RGB}{0,51,102}
\definecolor{darkred}{RGB}{180,30,30}
\definecolor{darkgreen}{RGB}{30,120,60}
\definecolor{darkorange}{RGB}{190,90,0}
\definecolor{lightgray}{RGB}{245,247,252}
\definecolor{warnyellow}{RGB}{255,243,205}
\definecolor{warnborder}{RGB}{200,160,0}

% ── Lệnh tiện ích ────────────────────────────────────────────
\newcounter{qnum}
\setcounter{qnum}{0}

% Hộp câu hỏi (màu đỏ nhạt)
\newcommand{\question}[2]{%
  \stepcounter{qnum}%
  \vspace{0.5em}
  \begin{mdframed}[backgroundcolor=darkred!8, linecolor=darkred,
    linewidth=1.8pt, leftmargin=0pt, rightmargin=0pt,
    innerleftmargin=12pt, innerrightmargin=12pt,
    innertopmargin=7pt, innerbottommargin=7pt]
  \noindent{\bfseries\color{darkred}[Câu \theqnum] #1}\\[2pt]
  {\small\itshape\color{darkblue}Loại câu hỏi: #2}
  \end{mdframed}
}

% Hộp trả lời (màu xanh nhạt)
\newcommand{\answer}[1]{%
  \begin{mdframed}[backgroundcolor=lightgray, linecolor=darkblue!30,
    linewidth=1pt, leftmargin=0pt, rightmargin=0pt,
    innerleftmargin=12pt, innerrightmargin=12pt,
    innertopmargin=8pt, innerbottommargin=8pt]
  \noindent{\bfseries\color{darkgreen}Gợi ý trả lời:}\\[4pt]
  #1
  \end{mdframed}
  \vspace{0.3em}
}

% Hộp cảnh báo / lưu ý chiến thuật
\newcommand{\tactic}[1]{%
  \begin{mdframed}[backgroundcolor=warnyellow, linecolor=warnborder,
    linewidth=1.2pt, leftmargin=0pt, rightmargin=0pt,
    innerleftmargin=10pt, innerrightmargin=10pt,
    innertopmargin=5pt, innerbottommargin=5pt]
  \noindent{\bfseries\color{darkorange}Chiến thuật trả lời:} \textit{#1}
  \end{mdframed}
  \vspace{0.4em}
}

% Tiêu đề phần
\newcommand{\qsection}[1]{%
  \vspace{1em}
  \begin{mdframed}[backgroundcolor=darkblue!90, linecolor=darkblue,
    leftmargin=0pt, rightmargin=0pt, innerleftmargin=10pt,
    innerrightmargin=10pt, innertopmargin=6pt, innerbottommargin=6pt]
  \color{white}\large\bfseries #1
  \end{mdframed}
  \vspace{0.3em}
}

% ─────────────────────────────────────────────────────────────
\begin{document}

\begin{center}
  {\LARGE\bfseries\color{darkblue}
    Chuẩn bị Hỏi đáp (Q\&A) --- Hội đồng Chấm điểm}\\[0.4em]
  {\large Đề tài: \textit{Một Khung Học Máy Tập Thể cho Ước Lượng}\\
    \textit{Công Sức Phần Mềm trên Ba Lược Đồ: LOC, FP, UCP}}\\[0.3em]
  {\normalsize Hội nghị Sinh viên Nghiên cứu Khoa học 2026 ---
    Trường Đại học Duy Tân}\\[0.15em]
  {\small\color{darkgreen}\bfseries Thời gian hỏi đáp: 3 phút
    (sau 2 phút demo) --- Tổng phần 2: 5 phút}
\end{center}
\hrule
\vspace{0.5em}

\begin{mdframed}[backgroundcolor=warnyellow, linecolor=warnborder,
  linewidth=1.5pt, innerleftmargin=12pt, innerrightmargin=12pt,
  innertopmargin=8pt, innerbottommargin=8pt]
\noindent{\bfseries\color{darkorange}Lưu ý chiến lược tổng quan:}
\begin{itemize}[leftmargin=1.5em, itemsep=1pt]
  \item Hội đồng xoáy vào: \textbf{(1) Chọn dataset}, \textbf{(2) Chọn thuật toán},
    \textbf{(3) Overfitting}, \textbf{(4) Cơ chế đánh giá mô hình}.
  \item Trả lời \textbf{ngắn gọn, thẳng vào kỹ thuật} --- tối đa 45 giây/câu.
  \item Phân công: người phụ trách Model/AI trả lời câu hỏi thuật toán;
    người phụ trách Data trả lời câu hỏi dataset.
  \item \textbf{Không} dùng từ ``em nghĩ'' --- dùng ``theo kết quả thực nghiệm của nhóm em\ldots''
\end{itemize}
\end{mdframed}

\vspace{0.5em}
\hrule
\vspace{0.3em}

% ═══════════════════════════════════════════════════════════════
\qsection{Nhóm 1 --- Câu hỏi về Dữ liệu và Dataset (G1)}
% ═══════════════════════════════════════════════════════════════

\question{Tiêu chí nào để nhóm chọn 18 nguồn dữ liệu này? Tại sao không dùng
  ít hơn hoặc nhiều hơn?}{Dataset Selection --- Mức độ khó: Trung bình}

\answer{
Nhóm em chọn theo ba tiêu chí cứng: \textbf{(1) có DOI hoặc URL có thể xác minh công khai},
\textbf{(2) chứa ít nhất một trong ba lược đồ LOC/FP/UCP}, và
\textbf{(3) không bị trùng lặp hoàn toàn với nguồn khác sau khi áp quy tắc deduplication}.
18 nguồn là tổng hợp của tất cả dataset phần mềm công khai thỏa ba tiêu chí trên tính đến
năm 2022, bao gồm các bộ nổi tiếng như ISBSG, NASA93, China, Desharnais.
Số lượng không cố định trước --- nhóm em dừng khi đã quét toàn bộ kho dữ liệu
PROMISE Repository và các bài báo trích dẫn từ năm 1979 đến 2022.
}

\tactic{Nhấn mạnh: ``Chúng em ghi lại toàn bộ nguồn gốc và viết rebuild script.
  Bất kỳ ai cũng có thể tái tạo chính xác bộ dữ liệu của chúng em từ đầu ---
  đây là Đóng góp số 1 của đề tài.''}

% ────────────────────────────────────────────────────────────────
\question{Nhóm xử lý trùng lặp dữ liệu giữa 18 nguồn như thế nào?
  Có rủi ro cùng một dự án xuất hiện ở nhiều nguồn không?}{Data Leakage --- Mức độ khó: Cao}

\answer{
Đây là rủi ro thực sự và nhóm em thừa nhận trong phần Threats to Validity.
Quy tắc deduplication của chúng em gồm: \textbf{(a) so khớp theo tên dự án + năm + effort},
\textbf{(b) loại bỏ nếu ba trường này trùng nhau ở ngưỡng $\pm$5\%},
và \textbf{(c) ưu tiên giữ bản từ nguồn có nhiều trường dữ liệu hơn}.
Tuy nhiên do thiếu metadata đầy đủ ở một số bộ cũ (NASA93, Kemerer), vẫn có khả năng
trùng lặp sót. Chúng em ghi nhận đây là giới hạn, và kịch bản rebuild script giúp
cộng đồng kiểm tra lại nếu muốn áp quy tắc nghiêm ngặt hơn.
}

% ────────────────────────────────────────────────────────────────
\question{FP chỉ có 158 dự án và UCP chỉ có 131 dự án --- liệu kết quả có đủ
  tin cậy về mặt thống kê không?}{Small Sample --- Mức độ khó: Cao}

\answer{
Nhóm em thiết kế giao thức riêng để giải quyết chính xác vấn đề này.
Với \textbf{FP} ($n=158$): dùng \textbf{Leave-One-Out Cross-Validation (LOOCV)}
--- 158 vòng, mỗi vòng test trên 1 điểm dữ liệu, cho phép tận dụng tối đa dữ liệu.
Kết quả FP được kèm \textbf{khoảng tin cậy 95\% bootstrap}.
Chúng em cũng ghi rõ kết quả FP là \textbf{mang tính tham khảo (exploratory)},
không khẳng định kết luận mạnh.
Với \textbf{UCP} ($n=131$): dùng 80/20 kết hợp 5-fold CV qua 10 seeds.
Đây là cách tiêu chuẩn cho tập dữ liệu nhỏ trong cộng đồng SEE.
}

\tactic{Nếu hội đồng hỏi thêm về mẫu nhỏ, hãy nói:
  ``Đây cũng là lý do chúng em chọn macro-averaging --- để không để LOC
  (2.765 dự án) che khuất kết quả thực sự của FP và UCP.''}

% ═══════════════════════════════════════════════════════════════
\qsection{Nhóm 2 --- Câu hỏi về Thuật toán và Kiến trúc (G2/G3)}
% ═══════════════════════════════════════════════════════════════

\question{Tại sao chọn Random Forest mà không dùng Deep Learning
  (LSTM, Transformer, hoặc Neural Network)?}{Algorithm Choice --- Mức độ khó: Cao}

\answer{
Ba lý do kỹ thuật cụ thể:
\begin{enumerate}[leftmargin=2em, itemsep=2pt]
  \item \textbf{Kích thước mẫu:} LSTM/Transformer cần hàng nghìn mẫu để học hiệu quả.
    FP và UCP chỉ có 131--158 dự án --- deep learning sẽ overfit nghiêm trọng.
  \item \textbf{Dữ liệu dạng bảng (tabular):} Theo nhiều benchmark tabular ML
    (Shwartz-Ziv \& Armon, 2022), Random Forest và XGBoost thường vượt deep learning
    trên dữ liệu bảng có ít đặc trưng ($< 20$ features).
  \item \textbf{Giải thích được:} Feature importance của RF giúp xác định biến nào
    đóng góp nhất --- quan trọng trong nghiên cứu khoa học SEE để cộng đồng kiểm chứng.
\end{enumerate}
Mục tiêu của nhóm không phải tìm mô hình mạnh nhất ---
mà xây dựng \textbf{khung so sánh công bằng và tái lập được}.
Deep Learning là hướng phát triển tương lai đã đề cập trong phần Future Work.
}

% ────────────────────────────────────────────────────────────────
\question{Tại sao không dùng COCOMO II đầy đủ làm baseline thay vì
  dùng power-law calibrated model?}{Baseline Design --- Mức độ khó: Cao}

\answer{
COCOMO II yêu cầu 17 cost driver (EM\textsubscript{i}) như kinh nghiệm đội, độ phức tạp nền tảng,
hỗ trợ công cụ, v.v. Tuy nhiên \textbf{phần lớn dataset FP và UCP công khai không có
các biến số này}. Nếu dùng giá trị mặc định ($EM_i = 1.0$), baseline sẽ vô tình yếu đi ---
đây là so sánh không công bằng, được gọi là ``straw-man comparison'' trong
cộng đồng nghiên cứu. Trong thí nghiệm sơ bộ, COCOMO II mặc định cho MMRE $> 2.5$
trên FP/UCP.

Thay vào đó, nhóm em dùng \textbf{power-law size-only model} được hiệu chỉnh
trên chính dữ liệu huấn luyện ($\log E = \alpha + \beta \log \text{Size}$)
--- bảo toàn tinh thần tham số của COCOMO nhưng công bằng hơn nhiều.
Đây là Đóng góp số 2 của đề tài.
}

% ────────────────────────────────────────────────────────────────
\question{Sự khác biệt giữa Random Forest và XGBoost trong bài toán này là gì?
  Tại sao RF lại tốt hơn XGBoost ở đây?}{Model Comparison --- Mức độ khó: Trung bình}

\answer{
Cả hai đều là ensemble methods nhưng khác nhau về cơ chế:
\textbf{Random Forest} xây song song nhiều cây và lấy trung bình (bagging) ---
ổn định hơn khi dữ liệu có nhiễu và mẫu nhỏ.
\textbf{XGBoost} xây tuần tự các cây bổ sung sai số (boosting) ---
mạnh hơn khi dữ liệu lớn và cần khai thác pattern tinh vi.

Trong thực nghiệm của nhóm em, RF đạt MAE $= 12.66$ PM so với XGBoost $= 13.21$ PM
trên macro-average. RF ổn định hơn (độ lệch chuẩn $\pm 0.85$ vs. $\pm 0.92$).
Điều này nhất quán với các nghiên cứu trước về SEE với tập dữ liệu không đồng nhất:
RF ít nhạy cảm hơn với ngoại lai và không cần tuning tinh vi để cho kết quả tốt.
}

% ────────────────────────────────────────────────────────────────
\question{Tại sao huấn luyện độc lập cho từng lược đồ mà không
  dùng transfer learning giữa LOC, FP và UCP?}{Cross-Schema Transfer --- Mức độ khó: Cao}

\answer{
Ba lược đồ có \textbf{không gian đặc trưng hoàn toàn khác nhau}:
LOC đo số dòng code (KLOC), FP đo điểm chức năng, UCP đo điểm use case ---
không có ánh xạ trực tiếp giữa các đơn vị này.
Nếu gộp chung và huấn luyện một mô hình duy nhất, mô hình sẽ học biểu diễn lẫn lộn
giữa ba hệ đo lường không tương thích.

Transfer learning giữa các lược đồ là \textbf{hướng nghiên cứu tương lai} thú vị ---
đã đề cập trong phần Future Work.
Để đảm bảo kết quả sạch cho benchmark này, nhóm em chọn
huấn luyện độc lập như một thiết kế thận trọng.
}

% ═══════════════════════════════════════════════════════════════
\qsection{Nhóm 3 --- Câu hỏi về Overfitting và Tổng quát hóa}
% ═══════════════════════════════════════════════════════════════

\question{Làm thế nào nhóm kiểm soát hiện tượng overfitting, đặc biệt
  với FP và UCP có mẫu rất nhỏ?}{Overfitting Control --- Mức độ khó: Cao}

\answer{
Nhóm em áp dụng \textbf{năm cơ chế kiểm soát overfitting đồng thời}:
\begin{enumerate}[leftmargin=2em, itemsep=2pt]
  \item \textbf{LOOCV cho FP ($n=158$):} 158-fold --- mỗi điểm được test độc lập.
  \item \textbf{10 random seeds độc lập:} Đảm bảo kết quả ổn định qua các phân chia khác nhau.
  \item \textbf{OOB score monitoring:} Out-of-Bag score của RF cho phép ước lượng
    lỗi trên dữ liệu chưa thấy mà không cần validation set riêng.
  \item \textbf{IQR-capping:} Giảm ảnh hưởng của ngoại lai cực đoan
    có thể khiến mô hình học các pattern bất thường.
  \item \textbf{Grid Search với inner CV (3-fold):} Hyperparameter tuning
    chỉ trên training fold, không để test set ảnh hưởng đến lựa chọn tham số.
\end{enumerate}
Mỗi cơ chế riêng lẻ có thể không đủ, nhưng kết hợp lại cho bằng chứng
vững chắc về tính ổn định của mô hình.
}

% ────────────────────────────────────────────────────────────────
\question{Leave-One-Source-Out (LOSO) là gì và tại sao nó quan trọng?
  Kết quả LOSO của nhóm nói lên điều gì?}{LOSO Generalization --- Mức độ khó: Trung bình}

\answer{
\textbf{LOSO} (Leave-One-Source-Out) là giao thức kiểm tra xuyên nguồn:
mô hình được huấn luyện trên dữ liệu từ 10 nguồn và kiểm tra trên
nguồn còn lại --- lặp lại 11 lần (vì LOC có 11 nguồn).
Đây là cách đo độ tổng quát hóa thực tế: liệu mô hình có thể dự đoán
tốt trên \textbf{dữ liệu từ tổ chức/dự án hoàn toàn mới} hay không.

Kết quả của nhóm: MAE tăng \textbf{chỉ 21\%} so với chia ngẫu nhiên 80/20.
Đây là mức chấp nhận được --- mô hình không bị phụ thuộc chặt vào đặc điểm
của từng nguồn cụ thể. Đây là câu trả lời cho \textbf{RQ2} của đề tài.
}

% ═══════════════════════════════════════════════════════════════
\qsection{Nhóm 4 --- Câu hỏi về Đánh giá và Số liệu}
% ═══════════════════════════════════════════════════════════════

\question{Tại sao báo cáo MMRE trong khi biết chỉ số này có bias
  (xu hướng ưu ái underestimation)?}{Evaluation Metric --- Mức độ khó: Trung bình}

\answer{
Nhóm em báo cáo MMRE chủ yếu vì \textbf{mục đích so sánh với tài liệu trước} ---
phần lớn các bài báo SEE dùng MMRE làm chỉ số chính, nên việc báo cáo giúp
cộng đồng so sánh trực tiếp.

Tuy nhiên, \textbf{chỉ số chính (primary metric) của nhóm em là MAE} ---
đối xứng, không bị lệch, và có ý nghĩa thực tiễn trực tiếp (person-months).
Nhóm em đã ghi nhận rõ giới hạn của MMRE trong phần Threats to Validity
(Construct Validity) và đề xuất sử dụng MdMRE cùng PRED(25) như chỉ số bổ sung
để giảm ảnh hưởng của bias.
}

% ────────────────────────────────────────────────────────────────
\question{Macro-averaging khác micro-averaging như thế nào?
  Tại sao macro-averaging công bằng hơn trong trường hợp này?}{Aggregation --- Mức độ khó: Trung bình}

\answer{
\textbf{Micro-averaging} tính trên toàn bộ 3.054 dự án gộp lại.
Do LOC chiếm 2.765/3.054 dự án ($= 90.5\%$), kết quả ``tổng thể'' thực chất
phản ánh gần như hoàn toàn kết quả trên LOC --- FP và UCP bị chìm.

\textbf{Macro-averaging} tính chỉ số riêng cho từng lược đồ rồi lấy trung bình đơn giản:
$$m_{\text{macro}} = \frac{1}{3}(m_{\text{LOC}} + m_{\text{FP}} + m_{\text{UCP}})$$
Mỗi lược đồ đóng góp $33.3\%$, bất kể kích thước mẫu.
Điều này đảm bảo đánh giá \textbf{công bằng ngang nhau} giữa các lược đồ ---
giải quyết trực tiếp Khoảng trống G3.
}

% ────────────────────────────────────────────────────────────────
\question{Kiểm định Wilcoxon được áp dụng như thế nào? Tại sao dùng
  Wilcoxon thay vì t-test thông thường?}{Statistical Test --- Mức độ khó: Trung bình}

\answer{
Nhóm em áp dụng \textbf{Wilcoxon signed-rank test} trên 10 cặp giá trị MAE
(RF vs. Baseline) từ 10 random seeds.
Wilcoxon được ưu tiên hơn t-test vì:
\textbf{(a)} Phân phối MAE qua các seeds không được đảm bảo là chuẩn (normal),
đặc biệt với tập dữ liệu nhỏ.
\textbf{(b)} Wilcoxon là non-parametric test --- không giả định phân phối chuẩn,
mạnh mẽ hơn khi $n = 10$ quan sát.

Kết quả: $p < 0.001$, $z = 9.87$ --- xác nhận sự cải thiện của RF so với baseline
\textbf{có ý nghĩa thống kê cao}. Hệ số hiệu ứng Cohen's $d = 1.23$ (lớn theo
ngưỡng của Cohen: $d > 0.8$) xác nhận kết quả không chỉ có ý nghĩa thống kê
mà còn có ý nghĩa thực tiễn.
}

% ────────────────────────────────────────────────────────────────
\question{PRED(25) là gì và tại sao chọn ngưỡng 25\%?}{Evaluation Metric --- Mức độ khó: Thấp}

\answer{
\textbf{PRED(25)} là tỷ lệ dự báo có sai số tương đối (MRE) $\leq 25\%$ so với thực tế:
$$\text{PRED}(k) = \frac{|\{i : \text{MRE}_i \leq k/100\}|}{n}$$
Ngưỡng 25\% là chuẩn ngành SEE từ nghiên cứu của Conte et al. (1986) và
được cộng đồng SEE duy trì để đảm bảo khả năng so sánh giữa các nghiên cứu.
Kết quả nhóm em: RF đạt PRED(25) $= 0.395$ --- nghĩa là $39.5\%$ dự án
được ước lượng sai không quá 25\%, so với $25.4\%$ của baseline.
}

% ═══════════════════════════════════════════════════════════════
\qsection{Nhóm 5 --- Câu hỏi về Đóng góp Khoa học và Ứng dụng}
% ═══════════════════════════════════════════════════════════════

\question{Đóng góp thực sự của đề tài so với các nghiên cứu SEE trước đây là gì?
  Điểm mới ở đâu?}{Research Contribution --- Mức độ khó: Cao}

\answer{
Nhóm em đóng góp \textbf{phương pháp luận} (methodological artifact),
không chỉ là một mô hình mới. Cụ thể, đây là \textbf{lần đầu tiên} một nghiên cứu
kết hợp đồng thời:
\begin{enumerate}[leftmargin=2em, itemsep=2pt]
  \item \textbf{Bản đồ nguồn dữ liệu có thể tái tạo} (dataset manifest + rebuild scripts)
    --- không chỉ ghi ``chúng tôi dùng PROMISE dataset''.
  \item \textbf{Baseline hiệu chỉnh công bằng} fitted trên training data per seed ---
    không dùng COCOMO II với tham số mặc định.
  \item \textbf{Giao thức riêng biệt cho từng lược đồ} (LOOCV cho FP nhỏ,
    LOSO cho LOC lớn).
  \item \textbf{Macro-averaging} để báo cáo công bằng qua ba lược đồ.
  \item \textbf{Tail evaluation} cho top 10\% dự án khó nhất.
\end{enumerate}
Khung này có thể tái sử dụng --- bất kỳ nghiên cứu nào muốn so sánh
mô hình mới đều có thể dùng framework của nhóm em làm nền tảng.
}

\tactic{Đây là câu quan trọng nhất. Hãy chủ động nhấn mạnh:
  ``Kết quả 42\% MMRE improvement là minh chứng cho giá trị của framework ---
  nhưng \emph{chính framework} mới là đóng góp cốt lõi.''}

% ────────────────────────────────────────────────────────────────
\question{Khung này có thể áp dụng cho dự án Agile/Scrum hiện đại không?
  Story points có thể tích hợp được không?}{Applicability --- Mức độ khó: Trung bình}

\answer{
Với \textbf{bộ dữ liệu hiện tại} ($1979$--$2022$, chủ yếu pre-Agile): không áp dụng trực tiếp.

Nhưng \textbf{thiết kế framework} hoàn toàn có thể mở rộng:
$s \in \{\text{LOC, FP, UCP, Story Points, Object Points}\}$ ---
chỉ cần thêm story points như lược đồ thứ tư với giao thức đánh giá riêng.
Thách thức là chưa có bộ dữ liệu Agile công khai đủ lớn và chuẩn hóa.
Đây là một trong những hướng nghiên cứu tương lai mà nhóm em đề xuất.
}

% ────────────────────────────────────────────────────────────────
\question{Kết quả thực nghiệm có bị ảnh hưởng bởi việc
  hầu hết dữ liệu đến từ Bắc Mỹ và Châu Âu không?}{External Validity --- Mức độ khó: Cao}

\answer{
Có --- đây là giới hạn rõ ràng mà nhóm em thừa nhận trong phần External Validity.
Phần lớn 18 nguồn đến từ Bắc Mỹ (NASA, IBM, Canada) và Châu Âu (Finland, UK),
ngoại trừ bộ China dataset.
Môi trường phát triển phần mềm tại Đông Nam Á hoặc các nước đang phát triển
có thể có đặc điểm năng suất, team size, và công nghệ khác.
Kết quả tổng quát hóa cho bối cảnh Việt Nam cần được kiểm chứng thêm
với dữ liệu địa phương --- đây cũng là hướng mà nhóm em muốn tiếp tục.
}

% ═══════════════════════════════════════════════════════════════
\qsection{Nhóm 6 --- Câu hỏi ``Bẫy'' / Bất ngờ}
% ═══════════════════════════════════════════════════════════════

\question{Mô hình của nhóm đạt $R^2 = 0.812$ --- nhưng $R^2$ cao có nghĩa là
  mô hình dự đoán tốt không? Có thể $R^2$ cao nhưng MAE vẫn lớn không?}{Critical Thinking --- Mức độ khó: Cao}

\answer{
Câu hỏi rất sắc bén. $R^2$ đo tỷ lệ phương sai được giải thích ---
$R^2$ cao nghĩa là mô hình bắt được xu hướng tổng thể, nhưng \textbf{không đảm bảo
sai số tuyệt đối nhỏ} trên từng dự án cụ thể.

Ví dụ: nếu dữ liệu có phân phối rất rộng (effort từ 1 đến 10.000 PM),
mô hình có thể đạt $R^2$ cao chỉ bằng cách đoán đúng chiều hướng,
trong khi MAE vẫn lớn về giá trị tuyệt đối.
Đây là lý do nhóm em \textbf{ưu tiên MAE và MdAE} làm chỉ số chính,
và báo cáo $R^2$ chỉ như chỉ số bổ sung.
}

% ────────────────────────────────────────────────────────────────
\question{Nếu nhóm em có thêm thời gian và nguồn lực, nhóm sẽ cải thiện điều gì
  trong nghiên cứu này?}{Future Work --- Mức độ khó: Thấp (nhưng quan trọng về ấn tượng)}

\answer{
Ba ưu tiên cụ thể và có thể thực hiện được:
\begin{enumerate}[leftmargin=2em, itemsep=2pt]
  \item \textbf{Transfer learning giữa các lược đồ:} Huấn luyện model trên LOC (lớn)
    rồi fine-tune cho FP/UCP (nhỏ) --- có thể tăng đáng kể hiệu năng trên FP/UCP.
  \item \textbf{Thu thập dữ liệu Agile:} Xây dataset story points từ Jira/Trello
    của các công ty phần mềm Việt Nam để kiểm chứng framework trong bối cảnh địa phương.
  \item \textbf{Tích hợp thêm đặc trưng:} Thêm team velocity, sprint length,
    technology stack --- thay vì chỉ dùng size metric.
\end{enumerate}
}

\tactic{Câu này là cơ hội thể hiện tư duy nghiên cứu. Trả lời tự nhiên,
  thể hiện nhóm hiểu giới hạn và có tầm nhìn rõ ràng.}

% ═══════════════════════════════════════════════════════════════
\vspace{1em}
\hrule
\vspace{0.5em}

\begin{mdframed}[backgroundcolor=darkblue!8, linecolor=darkblue,
  linewidth=1.5pt, innerleftmargin=12pt, innerrightmargin=12pt,
  innertopmargin=8pt, innerbottommargin=8pt]
\noindent{\bfseries\color{darkblue}Bảng tóm tắt phân công trả lời:}

\vspace{0.4em}
\begin{tabular}{>{\bfseries}p{4.5cm} p{9cm}}
\toprule
Loại câu hỏi & Người trả lời \\
\midrule
Dataset / Data quality & Thành viên phụ trách Data/Preprocessing \\
Thuật toán / Model & Thành viên phụ trách Model/AI \\
Evaluation metrics & Thành viên phụ trách Experiment/Evaluation \\
Đóng góp khoa học & Trưởng nhóm (người thuyết trình) \\
Ứng dụng thực tiễn & Bất kỳ thành viên nào --- chuẩn bị trước \\
\bottomrule
\end{tabular}
\end{mdframed}

\vspace{0.5em}

\begin{mdframed}[backgroundcolor=darkgreen!8, linecolor=darkgreen,
  linewidth=1.5pt, innerleftmargin=12pt, innerrightmargin=12pt,
  innertopmargin=8pt, innerbottommargin=8pt]
\noindent{\bfseries\color{darkgreen}Checklist cuối cùng trước hội nghị:}
\begin{itemize}[leftmargin=1.5em, itemsep=2pt]
  \item[$\square$] Tập trả lời từng câu theo đồng hồ --- mỗi câu $\leq 45$ giây.
  \item[$\square$] Mock presentation đủ 10 phút báo cáo + 2 phút demo.
  \item[$\square$] Chuẩn bị sẵn số liệu cụ thể: MAE $= 12.66$, MMRE $= 0.647$,
    $p < 0.001$, Cohen's $d = 1.23$, LOSO degradation $21\%$.
  \item[$\square$] Mọi thành viên thuộc lòng đóng góp chính của đề tài
    (5 đóng góp khoa học).
  \item[$\square$] Rebuild scripts và experiment logs sẵn sàng cung cấp cho hội đồng.
\end{itemize}
\end{mdframed}

\vspace{0.5em}
\hrule
\begin{center}
\textit{\small Tài liệu chuẩn bị Q\&A --- Hội nghị SVNCKH 2026 ---
  Trường Đại học Duy Tân, Tháng 4 năm 2026}
\end{center}

\end{document}
